\chap[implementation] Implementation

In this chapter, the technologies used in the implementation of NFC for the LionKey security key are described. 

\sec[nfc-specs] NFC technologies

After evaluating the available NFC operating modes and specifications, it becomes evident that the Card Emulation Mode is the only mode compatible with the intended use case. The objective is for the device to behave as a passive entity, analogous to a contactless smart card, rather than as an active initiator.

\nl In the ideal case, the use of Secure Element–based (SE) card emulation would be used. However, this is not feasible with respect to the existing LionKey hardware architecture. This approach requires the presence of a dedicated SE capable of isolated execution and secure key storage, which is not available in the current design. Consequently, the implementation must be based on host-based card emulation, where all protocol handling is performed by the microcontroller (MCU).

\nl The required functionality is achieved most appropriately through the emulation of a Type 4 NFC Forum Tag. In contrast to simpler tag types, Type 4 Tag operation is based on the ISO-DEP protocol and supports structured command–response communication via APDUs. This enables the implementation of a well-defined state machine and facilitates higher-level protocols, which are essential for authentication mechanisms such as FIDO2.

\nl Taking these constraints into account, the implementation adopts NFC-A as the underlying transmission technology. NFC-A provides the broadest interoperability across existing NFC readers, particularly in mobile devices, and is the de facto standard for ISO-DEP–based communication and for Type 4 tag implementations.

\sec[hardware] Hardware

The hardware chosen for the implementation is the STEVAL-D25R3916B, a daughter board of the STEVAL-25R3916B platform manufactured by STMicroelectronics \cite[STEVAL-25R3916B-brief]. This kit is an evaluation board for the 25R3916B NFC chip, which turned out to be the best option for the purpose of this thesis. 

\secc[board-specs] Board Specifications



\secc[alternatives] Alternative Chips

Although research on suitable hardware platforms for this project was conducted as part of the B4BPROJ6 course during the winter semester 2025/2026, a supplementary evaluation of alternative solutions was subsequently performed. This was primarily motivated by the inconsistencies identified in the software stack provided by STMicroelectronics, as will be discussed in Section \ref[middleware].

\nl Among the alternatives considered, the PN5180 was identified as a comparable solution, offering functionality equivalent \cite[PN5180-brief] to the selected ST-based hardware. A notable advantage of the NXP platform lies in the quality of its accompanying software libraries and reference implementations, which provide better support for development and integration. However, these are distributed under a license that restricts their use to NXP hardware, meaning that we would have to migrate the whole project to an NXP-manufactured MCU, which is beyond the scope of this bachelor thesis.

\nl Although alternative solutions are available, the selected approach prioritizes compatibility with existing hardware architecture and general modularity of the system. For this reason, hardware from STMicroelectronics was chosen, as it enables straightforward integration with the current platform on both the software and legal side.

\secc[note-energy-harvesting] Note: Energy Harvesting

\nl A significant limitation of the selected device is its absence of energy harvesting capabilities. Consequently, the security key cannot operate as a fully passive, standalone authenticator and requires an external power source during operation. It is important to note that achieving true standalone functionality would require not only an NFC frontend capable of efficient energy harvesting but also a substantially lower-power MCU. 

\nl The power demands of the current hardware platform exceed what can be reliably supplied by the electromagnetic field generated by typical NFC readers, making passive operation impossible under these conditions.
\sec[communication] Communication With the Board
\secc[nfc-spi] SPI
\secc[nfc-irq] IRQ
\sec[middleware] Middleware

